% s-309-p309-m001; original PDF 309, 310
\Aphorism{s-309-p309-m001}{362}
Lebensfeinde, asketische Priester, dithyrambische Künstler, deren Leidenschaft alles Leben überfluthet, Gefühlsekstasen, in denen das Leben nicht fester auf seiner Unterlage sitzt, als ein Mohnblatt an seinem Stengel, — das alles sind und waren immer noch die besten Verführer und Verführungen zum Leben! Wir Menschen sind wahrlich heroische Thiere: seitdem von allen Culturseiten \Sperr{gegen} das Leben geschrieben, gedacht, empfunden wird, können wir uns an ihm gar nicht satt leben. Je vorsorglicher man uns auf manches nicht Unbedenkliche hinweist und vor die gefährlichen Formen des Lebens warnende Vogelscheuchen hinpflanzt, desto fröhlicher klingt unsere Antwort: \Foreign{ergo vivamus}! Die theoretischen Vertheidiger des Lebens dürfen sich einer gleichen Wirksamkeit nicht rühmen, wohl aber eines gleichen Misserfolgs: seit die lauwarmen barmherzigen Seelen an der Vergutmüthigung und Weichpolsterung des Lebens arbeiten und alles Menschliche gern in einer grossen gut-geheizten Gotteskinderstube unterbringen möchten, seit man uns mit zuckersüssen Fadheiten den herben Wohlgeschmack des Lebens von der Zunge wegpinseln will, fängt das Leben an, wirklich und herzlich \Sperr{langweilig} zu werden, also unliebenswürdig, unlebenswürdig. Für diese Sippschaft ist, wie gesagt, der Mensch zu heroisch; ein Leben, wie sie es im Sinne haben, wohltemperirt, glatt, hedonistisch abgezirkelt und ausgerechnet, würde aufhören den „Willen zum Leben“ zu interessiren. Ich zweifle, ob irgend ein Rechtfertigungsphilosoph, ein Humanitätsmensch und Ausrufer des „stetigen Fortschritts“, ob irgend einer dieser weltbejahenden Teufelsadvocaten, die sich gern Optimisten nennen hören, uns zu seinem Teufel auch nur ein Fünkchen Lust und Liebe beigebracht hat; im Gegentheil, nichts discreditirt gründlicher als so eine Fürsprache. Aber den Pessimisten ist es geglückt, gerade ausgesuchte und genussfähige Geister zum Leben zu verlocken; ob sie nun Hadesschatten oder grelles quälendes Sonnenlicht auf das Dasein fallen liessen — jedenfalls gab es bei ihnen eine anregende Atmosphäre, nicht jenes schläfrige Regenwetter, in dem die optimistische Behaglichkeit mit Regenschirm und Stoicismus spazieren geht.
\par

